\documentclass[a4paper,landscape,twoside]{article}
\usepackage[margin=0mm]{geometry}
\usepackage{tikz}
\usepackage{graphicx}
\usepackage{fontspec}
\usepackage{polyglossia}
\setdefaultlanguage{czech}
\setmainfont{ {{- .BodyFontLatex -}} }[
  Ligatures      = TeX,
]
\setsansfont{ {{- .HeadingFontLatex -}} }[
  Ligatures      = TeX,
]
\usepackage[dvipsnames]{xcolor}
\usepackage{enumitem}
\usepackage{tabularx}
\usepackage{microtype}

% Bleed: extend paper by 3mm on each side with printer crop marks.
\usepackage[cam,noinfo,width=303truemm,height=216truemm,center]{crop}

% Broad printer compatibility
\pdfvariable minorversion 4

% Suppress timestamps for deterministic output
\pdfvariable suppressoptionalinfo 511

\pagestyle{empty}
\pagenumbering{gobble}

% \captionbadge — inline colored badge matching the photo overlay marker style.
% Renders a tight \colorbox with the marker number in bold contrast text,
% so footer caption markers visually match the on-photo overlay badges.
%   #1 = background color hex (without #)
%   #2 = text color name (white|black) — chosen via WCAG luminance in Go
%   #3 = marker string ("1", "1–3", "1, 3", etc.)
\newcommand{\captionbadge}[3]{%
  {\setlength{\fboxsep}{1.2pt}%
   \colorbox[HTML]{#1}{\textcolor{#2}{\textbf{#3}}}}%
}
% \captionbadgefb — fallback badge when no chapter color is set
% (matches the overlay fallback of black at 55% opacity).
%   #1 = marker string
\newcommand{\captionbadgefb}[1]{%
  {\setlength{\fboxsep}{1.2pt}%
   \colorbox{black!55}{\textcolor{white}{\textbf{#1}}}}%
}

\begin{document}
{{- range $si, $section := .Sections}}
{{- range $pi, $page := $section.Pages}}
% --- Page {{$page.PageNumber}} ({{$page.Style}}) ---
\begin{tikzpicture}[remember picture,overlay,shift={(current page.south west)},x=1mm,y=1mm]
% --- Canvas zone (clipped — expanded by heading bleed for colored boxes) ---
  \begin{scope}
    \clip ({{printf "%.2f" $page.ClipLeftX}},{{printf "%.2f" $page.CanvasBottomY}})
      rectangle ({{printf "%.2f" $page.ClipRightX}},{{printf "%.2f" $page.CanvasTopY}});
{{- range $si, $slot := $page.Slots}}
{{- if $slot.HasPhoto}}
{{- if $slot.IsArchival}}
  % Archival border
  \draw[black!45,line width=0.5pt]
    ({{printf "%.2f" $slot.BorderX}},{{printf "%.2f" $slot.BorderY}})
    rectangle ++({{printf "%.2f" $slot.BorderW}},{{printf "%.2f" $slot.BorderH}});
{{- end}}
  \begin{scope}
    \clip ({{printf "%.2f" $slot.ClipX}},{{printf "%.2f" $slot.ClipY}})
      rectangle ++({{printf "%.2f" $slot.ClipW}},{{printf "%.2f" $slot.ClipH}});
    \node[anchor=south west,inner sep=0]
      at ({{printf "%.2f" $slot.ImgX}},{{printf "%.2f" $slot.ImgY}})
      {\includegraphics[{{$slot.SizeDim}}={{printf "%.2f" $slot.SizeVal}}mm]{{print "{"}}{{$slot.FilePath}}{{print "}"}}};
  \end{scope}
{{- if gt $slot.CaptionMarker 0}}
  % Caption marker
{{- if $slot.ChapterColor}}
  \definecolor{markercolor}{HTML}{ {{- $slot.ChapterColor -}} }
  \fill[markercolor]
{{- else}}
  \fill[black,opacity=0.55]
{{- end}}
    ({{printf "%.2f" $slot.CaptionMarkerX}},{{printf "%.2f" $slot.CaptionMarkerY}})
    rectangle ++(4,4);
  \node[anchor=center,inner sep=0,font=\fontsize{6}{6}\selectfont\bfseries,text={{contrastTextColor $slot.ChapterColor}}]
    at ({{printf "%.2f" $slot.CaptionMarkerCenterX}},{{printf "%.2f" $slot.CaptionMarkerCenterY}})
    {{print "{"}}{{$slot.CaptionMarker}}{{print "}"}};
{{- end}}
{{- else if $slot.HasText}}
  % Text slot ({{$slot.TextType}})
  \node[anchor=north west,inner sep=0]
    at ({{printf "%.2f" (addFloat $slot.ClipX $slot.TextPadLeft)}},{{printf "%.2f" (addFloat $slot.ClipY $slot.ClipH)}})
    {\parbox{{print "{"}}{{printf "%.2f" (subtractFloat (subtractFloat $slot.ClipW $slot.TextPadLeft) $slot.TextPadRight)}}mm}{%
      \parskip=0pt\parindent=0pt\fontsize{ {{- printf "%.0f" $.BodyFontSize -}} }{ {{- printf "%.0f" $.BodyLineHeight -}} }\selectfont{{if $slot.RaggedRight}}\raggedright{{end}}%
      {{markdownToLatexColor $slot.TextContent $slot.ChapterColor $slot.BleedLeftMM $slot.BleedRightMM}}%
    }};
{{- end}}
{{- end}}
  \end{scope}
% --- Footer zone ---
{{- if $page.HasCaptions}}
  \node[anchor=north west,inner sep=0,text width={{printf "%.2f" $page.CaptionBlockW}}mm,
    font=\fontsize{ {{- printf "%.0f" $.CaptionFontSize -}} }{ {{- printf "%.0f" $.CaptionLeading -}} }\selectfont,text=black!{{$.CaptionOpacity}}]
    at ({{printf "%.2f" $page.CaptionBlockX}},{{printf "%.2f" $page.CaptionBlockY}})
    {{print "{"}}{{- range $ci, $cap := $page.Captions}}{{- if gt $ci 0}}\enspace {{end}}{{- if $cap.Marker}}{{- if $cap.ChapterColor}}\captionbadge{{print "{"}}{{$cap.ChapterColor}}{{print "}{"}}{{contrastTextColor $cap.ChapterColor}}{{print "}{"}}{{$cap.Marker}}{{print "}"}}{{- else}}\captionbadgefb{{print "{"}}{{$cap.Marker}}{{print "}"}}{{- end}}\, {{end}}{{latexEscape $cap.Caption}}{{- end}}{{print "}"}};
{{- end}}
  \node[anchor={{$page.FolioAnchor}},inner sep=0,font=\fontsize{9}{11}\selectfont,text=black!35]
    at ({{printf "%.2f" $page.FolioX}},{{printf "%.2f" $page.FolioY}})
    {{print "{"}}{{$page.PageNumber}}{{print "}"}};
{{- if $.DebugOverlay}}
% --- Debug overlay ---
  \draw[blue!40,line width=0.3pt]
    ({{printf "%.2f" $page.ContentLeftX}},{{printf "%.2f" $page.CanvasBottomY}})
    rectangle ({{printf "%.2f" $page.ContentRightX}},{{printf "%.2f" $page.CanvasTopY}});
  \draw[green!40,line width=0.2pt,dashed]
    ({{printf "%.2f" $page.ContentLeftX}},{{printf "%.2f" $page.HeaderY}})
    -- ({{printf "%.2f" $page.ContentRightX}},{{printf "%.2f" $page.HeaderY}});
  \draw[green!40,line width=0.2pt,dashed]
    ({{printf "%.2f" $page.ContentLeftX}},{{printf "%.2f" $page.FooterRuleY}})
    -- ({{printf "%.2f" $page.ContentRightX}},{{printf "%.2f" $page.FooterRuleY}});
{{- range $off := $.DebugColOffsets}}
  \draw[red!25,line width=0.2pt]
    ({{printf "%.2f" (addFloat $page.ContentLeftX $off)}},{{printf "%.2f" $page.CanvasBottomY}})
    -- ({{printf "%.2f" (addFloat $page.ContentLeftX $off)}},{{printf "%.2f" $page.CanvasTopY}});
{{- end}}
{{- end}}
\end{tikzpicture}
{{- if not $page.IsLast}}
\newpage
{{- end}}
{{- end}}
{{- end}}
\end{document}
